\chapter{Operations}
\label{ch:operations}

This chapter describes every operation. It groups the operations by what they
change. Each entry has the same points:

\begin{description}
\item[Writes] What the operation puts into the text.
\item[Refused] When it stops with an error and queues nothing. Chapter~\ref{ch:scripts}
  explains drift errors and refusals.
\item[was] What \code{was=} means, for the operations that have it.
\item[Again] What the operation does when you run it a second time on a design that
  already shows its result. It is \emph{skipped}, or it stops with a drift error.
\item[Example] The script and the transcript that show it.
\end{description}

\section{Rules that hold for all operations}

\begin{itemize}
\item An operation on a module handle changes every copy of the module. If the
  design instantiates the module more than once, you must confirm the number with
  \code{vet.module(name, instances=N)}. Without it, hdl-vet refuses the operation.
  The same holds for an operation on an instance handle. In this case, the module that
  holds the instance counts.
\item hdl-vet edits only text that it can edit safely. Chapter~\ref{ch:limits}
  collects these limits.
  \begin{itemize}
  \item hdl-vet refuses an operation on an instance when the instance comes from a
    macro, is written in a generate block, or is an instance array.
  \item It refuses an operation for connections when the instance uses positional or
    \code{.*} connections. It refuses an operation for overrides when the instance uses
    positional parameter overrides.
  \item It refuses an operation when the list that it edits holds a conditional
    compilation directive, for example \code{`ifdef}.
  \item It refuses an operation on a port list or a parameter list when the module
    header has a macro. Operations on ports need an ANSI port list.
  \end{itemize}
\item A file is a file on disk. A path is a place in the hierarchy (chapter~\ref{ch:intro}).
\item A name that does not exist raises an error. The error has up to three
  suggestions, in the form ``did you mean \ldots?''
\item An operation inserts text with the line ending of the file that the text goes
  into. hdl-vet never reformats text that it does not edit.
\item Each applied operation registers \emph{expectations} for the sign-off. These are
  the structural changes that the operation must cause. The entries name them as
  ``Sign-off requires'' and ``Sign-off permits''. Chapter~\ref{ch:signoff} explains
  how the sign-off uses them.
\end{itemize}

\section{All operations at a glance}

For each operation, the table shows three things:

\begin{itemize}
\item which text the operation edits
\item what \code{was} means for the operation, where it has one
\item what happens when you run the operation a second time on a result that already
  exists
\end{itemize}

The later sections give the details.

{\small
\begin{tabular}{@{}>{\raggedright\arraybackslash}p{0.305\textwidth}>{\raggedright\arraybackslash}p{0.2\textwidth}>{\raggedright\arraybackslash}p{0.28\textwidth}>{\raggedright\arraybackslash}p{0.12\textwidth}@{}}
\toprule
Operation & Edits the text of & \code{was} & Second run \\
\midrule
\apiname{Module.add_port} & the module & none & skipped \\
\apiname{Module.add_param} & the module & none & skipped \\
\apiname{Module.add_wire} & the module & none & skipped \\
\apiname{Module.add_assign} & the module & none & skipped \\
\apiname{Module.override_net} & the module & none & skipped, or drift \\
\apiname{Instance.connect} & the module that holds the instance & the expression now there. Without it, a different expression is drift & skipped \\
\apiname{Instance.disconnect} & the module that holds the instance & the expression now there, optional check & skipped \\
\apiname{Instance.set_param} & the module that holds the instance & the override now there. Without it, a different override is drift & skipped \\
\apiname{Vet.punch_port} & every module on the path & none & skipped \\
\apiname{Vet.propagate_param} & every module on the path above the destination & none & skipped \\
\apiname{Vet.add_file} & no source. The output tree gets a copy & none & skipped \\
\bottomrule
\end{tabular}

\medskip
\begin{tabular}{@{}>{\raggedright\arraybackslash}p{0.305\textwidth}>{\raggedright\arraybackslash}p{0.2\textwidth}>{\raggedright\arraybackslash}p{0.28\textwidth}>{\raggedright\arraybackslash}p{0.12\textwidth}@{}}
\toprule
Operation & Edits the text of & \code{was} & Second run \\
\midrule
\apiname{Instance.replace_module} & the module that holds the instance & the module now used, optional check & skipped, or drift \\
\apiname{Module.remove_instance} & the module that holds the instance & \code{module=}, the module now used, optional check & skipped \\
\apiname{Module.remove_wire} & the module & none & skipped \\
\apiname{Module.remove_assign} & the module & the right-hand side now there, optional check & skipped \\
\apiname{Module.remove_port} & the module, and every module that instantiates it & none & skipped \\
\apiname{Module.remove_param} & the module, and every module that instantiates it & none & skipped \\
\apiname{Vet.splice} & the file that you name & none. \code{old=} guards the bytes & none \\
\bottomrule
\end{tabular}
}

\section{Ports and connections}

\begin{entry}{\sig{Module.add_port}}
\item[Writes] The operation writes a port declaration at the end of the ANSI port list.
  It builds the declaration as direction, type, range and name, for example
  \code{input  logic [3:0] name}. The range comes from \code{width}:
  \begin{itemize}
  \item An integer greater than 1 gives \code{[N-1:0]}.
  \item A string \code{W} gives \code{[W-1:0]}.
  \item A string that starts with \code{[} gives that string as the range, as written.
  \item \code{None} or 1 gives no range.
  \end{itemize}
  If the closing parenthesis is alone at the start of its line, the new port
  goes on a new line, indented like the last port. A comma follows the last
  port. Otherwise the text \code{, declaration} follows the last port. An empty
  list gets the declaration between its parentheses.
\item[Refused] hdl-vet refuses the operation in these cases:
  \begin{itemize}
  \item a direction other than \code{input}, \code{output} or \code{inout}
  \item a module with a non-ANSI port list, or with none
  \item a header with a macro
  \item a conditional directive in the list
  \item a name that something else already uses, or a port of that name with a different
    declaration (a drift error)
  \item a name that is not a plain identifier, such as \code{"a b"}, or is a keyword,
    such as \code{"wire"}
  \item a name that begins with a backslash, or is not a \code{str}
  \end{itemize}
  The operation checks the name when you call it. It queues nothing for a name that it
  refuses.
\item[Again] hdl-vet skips the operation when the module has a port of that name with the
  same declaration.
\item[Example] \file{first.py} (chapter~\ref{ch:first}) and \file{ports.py} below.
  The operation does not connect a new input port in the instances of the module.
  Connect it with \apiname{Instance.connect}. If you do not, the sign-off reports a new
  diagnostic.
\end{entry}

\begin{entry}{\sig{Instance.connect}}
\item[Writes] The operation writes a connection \code{.port(expr)} after the last
  connection. It aligns the name field like the others when they all have the same width.
  An instance with no connections gets the connection between its parentheses. If the
  port is connected to nothing (\code{.port()}), the operation fills in the expression.
  If the port is connected to another expression, the operation replaces that
  expression, but only when \code{was} matches. Only the text between the parentheses
  changes.
\item[Refused] hdl-vet refuses the operation in these cases:
  \begin{itemize}
  \item a port that the module does not have
  \item an empty expression (to leave an output open, use \apiname{Instance.disconnect})
  \item the refusals for instances in the rules above
  \end{itemize}
\item[was] The expression that you expect to find now. When the port is connected to
  another expression, and you give no \code{was} or a different one, the operation stops
  with a drift error. The error quotes what it found. The operation fills an open port
  (\code{.port()}) without \code{was}.
\item[Again] hdl-vet skips the operation when the port already has the same expression.
\item[Example] \file{first.py}, and \file{drift.py} in chapter~\ref{ch:scripts}.
\end{entry}

\begin{entry}{\sig{Instance.disconnect}}
\item[Writes] The operation writes an empty connection \code{.port()}. The connection
  keeps the name field.
\item[Refused] The operation refuses an input port, because an open input floats. The
  message names the instance and the port. It gives two other choices:
  \begin{itemize}
  \item tie the input to a value with \apiname{Instance.connect}
  \item remove the port from the module
  \end{itemize}
  The operation refuses an interface port too.
\item[was] The expression that you expect to find now. A different expression is a drift
  error. Without \code{was}, the operation does not check. It opens the port whatever is
  connected.
\item[Again] hdl-vet skips the operation when the port is already open or has no connection.
\item[Example] \file{ports.py}. After the change, \code{count1_o} of the top
  module has no driver. This is allowed, because a top output that nobody reads can float.
  If something reads the signal, the sign-off says \emph{left dangling}
  (chapter~\ref{ch:signoff}).
\end{entry}

\lstinputlisting[style=script]{examples/scripts/ports.py}

\section{Parameters}

\begin{entry}{\sig{Module.add_param}}
\item[Writes] The operation writes \code{parameter type name = default} after the last
  parameter in the parameter port list. If the module has no such list, the operation
  writes \texttt{ \#(parameter type name = default)} directly after the name of the
  module, as in the transcript below. \code{default} can be any value. The operation
  writes it as text.
\item[Refused] hdl-vet refuses the operation in these cases:
  \begin{itemize}
  \item a header with a macro
  \item a conditional directive in the list
  \item a name that another parameter in the body, or something else, already uses (a
    drift error)
  \item a name that is not a plain identifier, such as \code{"a b"}, or is a keyword,
    such as \code{"module"}
  \item a name that begins with a backslash, or is not a \code{str}
  \end{itemize}
  The operation checks the name when you call it. It queues nothing for a name that it
  refuses.
\item[Again] hdl-vet skips the operation when the module has a parameter of that name with
  the same type and default.
\item[Example] \file{params.py}.
\end{entry}

\begin{entry}{\sig{Instance.set_param}}
\item[Writes] The operation writes an override \code{.name(value)} after the last
  override. When the instance has no overrides, it writes \texttt{ \#(.name(value))}
  directly after the module name. It replaces an existing override when \code{was}
  matches.
\item[Refused] hdl-vet refuses the operation in these cases:
  \begin{itemize}
  \item a parameter that the instantiated module does not have, a \code{localparam}, or a
    type parameter
  \item positional overrides
  \item an instance that shares its parameter list with other instances in one statement
  \item the refusals for instances in the rules above
  \end{itemize}
\item[was] The override text that you expect to find now. A different text is a
  drift error.
\item[Again] hdl-vet skips the operation when the override already has the same text.
\item[Example] \file{params.py}. The sign-off compares the text of the override,
  not the resolved value. A change of the parameter can change port widths below
  the instance. Declare these changes with \apiname{Vet.expect} (see
  section~\ref{sec:byhand}).
\end{entry}

\lstinputlisting[style=script]{examples/scripts/params.py}

\section{Wires and assignments}

\begin{entry}{\sig{Module.add_wire}}
\item[Writes] The operation writes a declaration such as \code{logic [3:0] name;} on a
  new line directly below the line that ends the module header. It indents the
  declaration like the first line of the body. If the body is empty, it indents with two
  blanks.
\item[Refused] hdl-vet refuses the operation in these cases:
  \begin{itemize}
  \item a name that is not a plain identifier, or is a keyword, or begins with a
    backslash, or is not a \code{str}
  \item a module header whose last line holds the beginning of the body, or ends inside a
    block comment
  \item a name that something else uses, or a net of that name with another type (a drift
    error)
  \end{itemize}
\item[Again] hdl-vet skips the operation when the module has a net or variable of that name
  with the same type and range.
\item[Example] \file{skipped.py} (chapter~\ref{ch:scripts}).
\end{entry}

\begin{entry}{\sig{Module.add_assign}}
\item[Writes] The operation writes \code{assign target = expr;} on a new line directly
  above the line of \code{endmodule}.
\item[Refused] hdl-vet refuses the operation in these cases:
  \begin{itemize}
  \item a target that is not a net, variable or port of the module (with suggestions)
  \item an input port
  \item a target that already has a driver
  \item a module whose \code{endmodule} does not start its line
  \end{itemize}
  The target must be a plain identifier.
\item[Again] hdl-vet skips the operation when the module has exactly that assignment.
\item[Example] \file{ports.py}.
\end{entry}

\section{Logic}
\label{sec:logic}

The operation in this section changes what the design computes. The sign-off proves that
the change has no effect when you switch it off (chapter~\ref{ch:signoff}).

\begin{entry}{\sig{Module.override_net}}
\item[Writes] The operation puts a multiplexer on a net. The net keeps its name and its
  readers. A new signal, \code{keep}, holds the value that the driver gives. The default
  name of \code{keep} is the name of the net with \code{_pre} at the end. The operation
  makes three changes in the text of the module:
  \begin{itemize}
  \item It declares \code{keep} with the type and the width of the net. The declaration is
    on a new line below the declaration of the net. For a port, it is below the line
    that ends the module header.
  \item It makes the driver of the net drive \code{keep}. The driver is an \code{assign}
    statement, or the connection of an output pin of an instance.
  \item It writes \code{assign net = enable ? value : keep;} on a new line below the
    driver.
  \end{itemize}
  \code{value} is the text of an expression over signals of the module. \code{enable} is
  the name of an input port of the module. The port is one bit wide. Add it first, with
  \apiname{Module.add_port} or \apiname{Vet.punch_port}. When \code{enable} is 0, the net
  has its old value. When \code{enable} is 1, the net has \code{value}.

  The input design must not have the port of \code{enable}. The proof ties the enable to 0
  in the edited module only.

  \code{value} is one expression, and it has the width of the net. A fill literal
  (\code{'0}, \code{'1}, \code{'x} or \code{'z}) takes the width of the net. hdl-vet finds
  the width of the value in the compiler. It does not guess from the text.
\item[Refused] hdl-vet refuses the operation in these cases:
  \begin{itemize}
  \item a net that the module does not have, or an input port
  \item a net with no driver, with more than one driver, or with a driver for a part of
    the net only
  \item a net that an \code{always} block drives
  \item a net that an inout pin of an instance connects
  \item a net that is an unpacked array, or has an initial value
  \item a net that has an override already. This is a drift error. A net gets one override.
  \item a driver in a generate block or in an instance array
  \item a net that has an unpacked dimension or an initial value
  \item an \code{enable} that is not an input port of the module, or is wider than one bit
  \item an \code{enable} that the input design has as a port already
  \item a \code{keep} that is not a plain identifier, or that the module uses already
  \item an empty \code{value}
  \item a \code{value} that is not one expression. Examples are a comma at the top level, an
    unbalanced bracket, a second statement and a comment.
  \item a \code{value} with a line break. A value is one line.
  \item a \code{value} that reads the net itself, or a name the net drives inside the
    module. The message says \emph{loop}. hdl-vet reads the names from the syntax of the
    value, not from its text. It judges the value by the name each path starts with. A
    route through a child instance, a hierarchical name or a function call is not seen
    (chapter~\ref{ch:limits}).
  \item a \code{value} with a width other than the width of the net. The message says
    \emph{width}.
  \item a module that the design instantiates more than once, when you did not give the
    count (see the rules above)
  \item a \code{waiver} that is not a text, or has a line break
  \item a call without a waiver when Yosys is not installed. The message names
    \code{hdl-vet[equiv]}.
  \end{itemize}
  The operation checks all this when you call it. It queues nothing for a call that it
  refuses. Like the removals, it loads the edited design. If the design has a new
  diagnostic, the operation takes its edits back and stops.
\item[Sign-off requires] When an instance pin drives the net, the sign-off requires a
  \code{connection-changed} of that pin at every path of the module. An \code{assign}
  driver needs no change of the structure.
\item[Waiver] \code{waiver} is a text that says why the proof is not needed, for example
  the name of a test that covers the change. hdl-vet prints it in the report and writes it
  to \file{hdl-vet.json}. With Yosys, the check still runs and the report shows what it said. Only the waiver
  decides whether the sign-off passes. Without Yosys, the operation needs a waiver.
  Chapter~\ref{ch:signoff} gives the rules.

  The waiver is one line. The report shows it in quotes, as JSON does. The waiver of the
  latest call counts. A call that is skipped sets the waiver of that
  override to what the call gives. A call without a waiver takes the waiver away.
\item[The way of the enable] The proof ties the enable to 0. The sign-off also follows the
  enable up to the top module (section~\ref{sec:enable}). It looks at the pin of the enable
  on each instance of the module in every elaboration of the edited design. A pin is
  \emph{off} when it connects a literal with the value 0. A pin is \emph{carried} when the
  compiler resolves its connection to a new input port of the parent. Any other connection is \emph{not off}, and the sign-off fails unless every
  operation that uses the enable has a waiver.
\item[was] The operation has no \code{was}.
\item[Again] hdl-vet skips the operation when the net is driven by the line
  \code{assign net = enable ? value : keep;} that the operation writes. If that line has
  another enable, another value or another kept signal, the operation stops with a drift
  error.
\item[Example] \file{override.py} bypasses the reset of the example design. The script
  \file{override_again.py} runs after it and is skipped. The script
  \file{override_waiver.py} gives a waiver. The script \file{override_refused.py} names an
  enable with two bits. The script \file{override_wrong.py} swaps the two branches of the
  multiplexer by hand. The sign-off then says that the module is not equivalent.

  The script \file{override_levels.py} puts one enable on an override in \code{rst_sync} and
  on an override in \code{soc_top}. The script \file{override_not_off.py} connects the pin of
  the enable of \code{rst_sync} to 1. The sign-off then says that the enable is not off.
\end{entry}

\lstinputlisting[style=script]{examples/scripts/override.py}
\lstinputlisting[style=script]{examples/scripts/override_again.py}
\lstinputlisting[style=script]{examples/scripts/override_waiver.py}
\lstinputlisting[style=script]{examples/scripts/override_refused.py}
\lstinputlisting[style=script]{examples/scripts/override_wrong.py}
\lstinputlisting[style=script]{examples/scripts/override_levels.py}
\lstinputlisting[style=script]{examples/scripts/override_not_off.py}

\lstinputlisting[style=console]{examples/console/override.txt}

The first command shows the diff and the line of the check. The second command shows the
second run. The third command shows a waiver. The fourth command shows a refusal. The fifth
command shows a failed check. The sixth command shows two levels, and the seventh command
shows an enable that is not off.

\section{Through the hierarchy}

A generator makes the hierarchy. A change deep inside the hierarchy means a port or a
parameter in every module on the way. The next two operations do all of this.

The two operations take a hierarchical path \code{to}. The path starts with the name of
the top module, for example \code{soc_top.u_periph.u_timer0}. The modules on the path are
$M_0$ (the top module) to $M_n$ (the destination). hdl-vet checks the whole path
before it changes anything.

\begin{entry}{\sig{Vet.punch_port}}
\item[Writes] For an \emph{input}, a signal enters at the top module and arrives at
  $M_n$.

  $M_0$ gets the input port \code{name}, unless you give \code{signal}. \code{signal} is
  the name of an existing port or net of $M_0$ that is the source. In this case $M_0$
  gets no new port. Each of $M_1\ldots M_n$ gets the input port \code{name}. The
  operation connects the instance of each of these modules in the module above to
  \code{name}. If you gave \code{signal}, it connects the instance in $M_0$ to
  \code{signal}.

  For an \emph{output}, \code{signal} (default \code{name}) is a port or net of $M_n$.
  $M_n$ gets the output port \code{name} and \code{assign name = signal;}. Each of
  $M_{n-1}\ldots M_0$ gets the output port. The operation connects the instance below to
  it.

  If you give \code{type} or \code{width}, the operation declares the port as
  \apiname{Module.add_port} does. If you give neither and there is a signal, the port
  takes the type of that signal as written (an output always has one). Without a signal
  the input is a 1-bit \code{logic}.

  The type must mean the same in every module on the path. For this, you must write
  every identifier in the type as \code{pkg::NAME}. The operation refuses a type that
  uses a parameter or typedef of one module. In this case, give \code{width=}.
\item[Refused] hdl-vet refuses the operation in these cases:
  \begin{itemize}
  \item a path that is the top alone, does not start at the top module, or names an
    instance that does not exist
  \item a segment with an index, instances in generate blocks and instance arrays
  \item an interface port
  \item a signal with unpacked dimensions
  \item a module on the path that the design instantiates more than once
  \end{itemize}
  For such a module, the operation refuses an input unless you give
  \code{all_instances=True}. It always refuses an output, because several copies then
  drive one net. With \code{all_instances=True}, the operation connects every instance of
  each module on the path. Each of these instances must be in a module of the path.
\item[Again] A second identical call skips every piece. If a path is partly
  punched, the operation completes it. If the output signal has the same name as an
  existing output port of $M_n$, the operation does not edit $M_n$.
\item[Example] \file{punch.py} and \file{punch_refused.py}.
\end{entry}

\lstinputlisting[style=script]{examples/scripts/punch.py}

\lstinputlisting[style=script]{examples/scripts/punch_refused.py}

\lstinputlisting[style=console]{examples/console/hierarchy.txt}

In the transcript, \file{punch_refused.py} catches the two refusals and prints
them. The design instantiates the module \code{timer} twice. The input needs
\code{all_instances=True}. The output cannot pass through two copies at all.

\file{punch.py} then does what is possible. It brings the signal
\code{rst_no} out of \code{u_rst_sync} to a new output of the top module. It
adds an input to every copy of \code{timer} down from the top. In the listing of
\file{soc_top.sv}, look at the new lines:

\begin{itemize}
\item the output port \code{sync_rst_o} and the input port \code{tick_i} at the end of
  the port list
\item the connection \code{.sync_rst_o(sync_rst_o)} on \code{u_rst_sync}
\item the connection \code{.tick_i(tick_i)} on \code{u_periph}
\end{itemize}

\begin{entry}{\sig{Vet.propagate_param}}
\item[Writes] The operation makes a parameter of $M_n$ settable from the top, without a
  change of its value. Each of $M_0\ldots M_{n-1}$ gets
  \code{parameter type name = default}. The operation overrides the instance below with
  \code{.name(name)}. The default is the first of these that exists:
  \begin{enumerate}
  \item \code{default}, if you give it
  \item the value that overrides the destination instance now
  \item the default that $M_n$ declares
  \end{enumerate}
  The type is the type of $M_n$. An implicit type stays implicit.
\item[Refused] hdl-vet refuses the operation in these cases:
  \begin{itemize}
  \item a parameter that $M_n$ does not have, a \code{localparam}, a type
    parameter, or a parameter declared in the body
  \item a type or a value that is not the same in every module on the path (the rules of
    \apiname{Vet.punch_port}). In this case, give \code{default=}.
  \end{itemize}
  The design must instantiate each of $M_0\ldots M_{n-1}$ once. The design can
  instantiate $M_n$ itself more than once. The operation does not edit it.
\item[Again] hdl-vet skips the operation when the path already passes the parameter down.
\item[Example] \file{propagate.py}. The value in use must not change. If it changes,
  the sign-off says so. The sign-off does not permit a parameter change at the
  destination, and a check compares the resolved value before and after.
\end{entry}

\lstinputlisting[style=script]{examples/scripts/propagate.py}

\lstinputlisting[style=console]{examples/console/ports.txt}

In the transcript, look at three diffs:

\begin{itemize}
\item \file{ports.py} adds the port \code{status_o} with its assignment, and opens
  \code{count1_o} on \code{u_periph}.
\item \file{params.py} turns the module header into
  \code{module soc_top }\texttt{\#}\code{(parameter int DEPTH = 2) (} and uses the
  parameter in the override.
\item \file{propagate.py} adds a parameter \code{WIDTH} to \code{periph_wrap} and to
  \code{soc_top}, and passes it down. In \code{soc_top} the instance becomes
  \code{periph_wrap }\texttt{\#}\code{(.WIDTH(WIDTH)) u_periph}. In \code{periph_wrap}
  the first timer gets \code{.WIDTH(WIDTH)}. The value stays 8 everywhere.
\end{itemize}

\section{Replacing a module}

You have a module that is not in your design, a technology cell. You want an
instance to use it. The cell is in a file that is not in the file list of the
generator. Two operations do the job. \apiname{Vet.add_file} adds the file to the
design. \apiname{Instance.replace_module} changes the module of the instance.

\lstinputlisting[style=source]{examples/cells/tech_rst_sync.sv}

\begin{entry}{\sig{Vet.add_file}}
\item[Writes] The operation writes nothing in the sources. The file joins the design as
  its last source. hdl-vet never edits it. These things happen in the output tree:
  \begin{itemize}
  \item The output tree has a copy of the file at \file{hdl_vet_added/NAME}.
  \item The file list gets one more line for the file, after the files that belong to the
    design.
  \item \file{hdl-vet.json} lists the file under \code{added}.
  \end{itemize}
  On the console, hdl-vet prints \code{added hdl_vet_added/NAME}. A relative file name
  in the argument \code{path} starts from the folder of the script. It starts from the
  current folder when no file of that name is next to the script. An absolute file name
  stays as it is. Outside a script, a relative file name starts from the current folder.
\item[Refused] hdl-vet refuses the operation in these cases:
  \begin{itemize}
  \item a relative file name that is next to the script and also in the current folder, but
    not one file: give an absolute file name
  \item a file that cannot be read
  \item a second added file with the same name
  \item an input tree that has the folder \file{hdl_vet_added} at its root
  \item a file that includes a header that is not a file of the design
  \item a file whose addition gives an error: a module that exists already, or a syntax
    error
  \end{itemize}
  hdl-vet does not count a warning as new if the file had it when the script added the
  file. hdl-vet counts a warning that appears in the file later, for example after an
  edit elsewhere changes how the file is read.
\item[Again] hdl-vet skips the operation when the file is already part of the design.
\item[Sign-off] The sign-off fails with \code{failed check: add_file NAME: none of
  its modules is instantiated} when no instance, other than a top module, uses a module
  from the file. The last transcript below shows this failure.
\item[Example] \file{swap_cell.py}, \file{add_cell.py}.
\end{entry}

\begin{entry}{\sig{Instance.replace_module}}
\item[Writes] The operation writes the name of the new module in place of the old name.
  It also writes the changes to the connections and overrides that you name. The steps run
  in this order:
  \begin{enumerate}
  \item The operation replaces the module name.
  \item It gives a new name to renamed connections and overrides. The blanks before the
    opening parenthesis stay.
  \item It removes dropped connections and overrides. If an override list becomes empty,
    it removes the whole list.
  \item It adds \code{set_params}, then \code{connect}, then \code{unconnected}. It adds
    the entries of each in the order given.
  \end{enumerate}
  The operation does not reformat existing text.
\item[Rules] Each connection and each override of the instance, and each port
  of the new module, must be covered by one of these rules:
  \begin{itemize}
  \item \code{ports={"old": "new"}} renames a connection. \code{drop_ports=["old"]}
    removes it. The operation keeps a connection that you name in neither, if the new
    module has a port with that name.
  \item \code{params} and \code{drop_params} do the same for parameter overrides.
    This holds for an override of a type parameter too. The operation keeps it when the
    new module has a type parameter of that name. \code{params} renames it onto a type
    parameter. The text of the override does not change.
  \item Every port of the new module needs a connection. The connection is one of these:
    \begin{itemize}
    \item the target of a kept or renamed connection
    \item \code{connect={"new": expr}}
    \item \code{unconnected=["new"]}
    \end{itemize}
    \code{unconnected} is for outputs and inouts. It writes an empty connection.
  \item \code{set_params={"name": value}} adds an override, or replaces a kept one.
    Parameters of the new module that nothing overrides keep their defaults.
  \end{itemize}
\item[Refused] hdl-vet checks the maps first. It reports all problems in one
  message. The problems are:
  \begin{itemize}
  \item something not covered by a rule (the message says \emph{is not accounted for} and
    names each one)
  \item a name that the instance or the new module does not have
  \item an input in \code{unconnected}
  \item two things on one port or parameter
  \item a \code{set_params} name that is not an overridable value parameter
  \item a \code{params} map from a type parameter to a value parameter, or the reverse
  \end{itemize}
  The operation also refuses a module with a non-ANSI port list, an interface port, and
  positional overrides. It tries the edit. It refuses the edit when a kept or renamed port
  has another direction, width or resolved type in the new module. A typedef or a type
  parameter with the same resolved type is no difference. The message names both ports.
  It also refuses the edit when the edit gives another value to a parameter that both modules
  have and nothing overrides.
\item[accept\_types] A list of ports of the new module. A renamed port keeps the name it
  lands on. The list waives a difference of type for these ports only. A difference of
  direction or bit width is never waived. Each waived port adds a note with both types to the
  sign-off report. The operation refuses a name that is not a port of the new module, and a
  waiver that covers nothing.
\item[was] The name of the module that you expect the instance to use now. If the
  instance uses neither the module in \code{was} nor the new module, the operation
  stops with a drift error. The error names the module that the instance uses. Without
  \code{was} the operation does not check. It replaces whatever module the instance uses.
\item[Again] When the instance already uses the new module, hdl-vet compares only
  \code{connect} and \code{set_params}. If every entry is there, hdl-vet skips the operation.
  Otherwise it stops with a drift error that names the first entry that differs or is
  missing.
\item[Sign-off] The sign-off requires \code{module-changed} at every path of the
  instance. It permits every change at these paths and below them. At the end, a check
  confirms that the instance uses the new module.
\item[Example] \file{swap_cell.py}, \file{swap_cell_drop.py} and
  \file{swap_wrong.py}.
\end{entry}

\lstinputlisting[style=script]{examples/scripts/swap_cell.py}

\lstinputlisting[style=script]{examples/scripts/swap_cell_drop.py}

\lstinputlisting[style=script]{examples/scripts/swap_wrong.py}

\lstinputlisting[style=script]{examples/scripts/add_cell.py}

\lstinputlisting[style=console]{examples/console/replace.txt}

In the first script, the ports \code{clk_i}, \code{rst_ni}, \code{test_i} and
\code{rst_no} get the names that the cell uses, and the parameter \code{STAGES}
becomes \code{DEPTH}. The cell has an extra output, \code{SO}. The script
leaves it open with \code{unconnected}.

The second script does the same in another way. It drops \code{test_i}, and it ties
\code{SE} of the cell to a constant with \code{connect}. It also drops the parameter
\code{STAGES} and sets \code{DEPTH}.

In the transcript, you see three things:

\begin{itemize}
\item the diff of the first script
\item the output tree of the second script
\item two failures: \file{swap_wrong.py} does not cover \code{SO} and
  \code{STAGES}. \file{add_cell.py} adds the cell and does not use it.
\end{itemize}

Look at the diff of the first script. The renamed connections, such as \code{.CK  (clk_i)},
keep the blanks that the old names had, so they are not aligned with each other any more.
This is as expected. hdl-vet never reformats bytes that it does not edit, and a rename
changes only the name. The line \code{added hdl_vet_added/tech_rst_sync.sv} after the
diff says that the cell goes into the output tree.

\section{Removing things}

The removal operations are checked operations (chapter~\ref{ch:scripts}). hdl-vet tries
each one on the design. If the operation leaves something dangling, gives a new
diagnostic, or leaves a signal partly undriven, hdl-vet refuses it. It takes back
everything that the operation did. To remove a driver and its reader together, use a group.

The removal takes the lines that it empties. This is the \emph{line rule}. If only
blanks stay on the lines that the removed text occupied, the removal takes these lines
too, with one line break. For an item of a list (a port, parameter, connection or
override), the rule is:

\begin{itemize}
\item If the item is not the last, the removal takes the text from the start of the item
  to the start of the next item.
\item If the item is the last of several, the removal takes the comma after the previous
  item, with the item.
\item If the item is the only item, the removal takes the item. If a list of overrides
  becomes empty, the removal takes the whole list. If a parameter port list becomes empty
  and only blanks stand between the name and \texttt{\#}, the removal takes the whole
  list. This is the text from the end of the module name through the closing parenthesis.
\end{itemize}

\begin{entry}{\sig{Module.remove_instance}}
\item[Writes] The operation removes the instantiation statement (line rule).
\item[Refused] hdl-vet refuses the operation in these cases:
  \begin{itemize}
  \item an instance that comes from a macro, a generate block or an array
  \item a statement that declares several instances
  \end{itemize}
\item[was] \code{module} works as \code{was}. If you give it and the instance uses
  another module, the operation stops with a drift error. Without it, the operation does
  not check which module the instance uses.
\item[Again] hdl-vet skips the operation when the module has no instance of that name.
\item[Sign-off] The sign-off requires \code{instance-removed} at every path. It permits
  any change below these paths.
\item[Example] \file{group.py}.
\end{entry}

\begin{entry}{\sig{Module.remove_wire}}
\item[Writes] The operation removes the declaration of a module-level net or variable
  (line rule).
\item[Refused] hdl-vet refuses the operation in these cases:
  \begin{itemize}
  \item a name that is a port or something other than a net or variable
  \item a declaration that declares several names, comes from a macro or an include, or
    is in a generate block
  \end{itemize}
\item[Again] hdl-vet skips the operation when nothing of that name is declared.
\item[Example] \file{group.py}.
\end{entry}

\begin{entry}{\sig{Module.remove_assign}}
\item[Writes] The operation removes the continuous assignment to \code{target} (line
  rule).
\item[Refused] hdl-vet refuses the operation in these cases:
  \begin{itemize}
  \item several assignments to the target
  \item a statement that holds several assignments, comes from a macro or an include, or
    is in a generate block
  \end{itemize}
\item[was] The right-hand side that you expect. A different one is a drift error.
  Without \code{was}, the operation does not check.
\item[Again] hdl-vet skips the operation when there is no assignment to the target.
\item[Example] \file{retire.py}.
\end{entry}

\begin{entry}{\sig{Module.remove_port}}
\item[Writes] The operation removes the port from the ANSI port list, and the connection
  to it at every instance of the module.
\item[Refused] hdl-vet refuses the operation in these cases:
  \begin{itemize}
  \item a port whose successor takes its direction or type from it (a declaration like
    \code{output logic [3:0] y0, y1})
  \item a non-ANSI port list
  \item a header with a macro
  \item any instance whose connection list cannot be edited
  \end{itemize}
\item[Again] hdl-vet skips the operation when the module has no such port.
\item[Sign-off] The sign-off requires \code{port-removed} at every path when the input
  design had the port. It permits \code{connection-changed} for the port. When
  \apiname{Module.add_port} of this run added the port, the removal forgets that
  expectation instead, so the pair leaves nothing behind.
\item[Example] \file{retire.py}.
\end{entry}

\begin{entry}{\sig{Module.remove_param}}
\item[Writes] The operation removes the parameter from the parameter port list, and its
  override at every instance.
\item[Refused] hdl-vet refuses the operation in these cases:
  \begin{itemize}
  \item a parameter that is declared in the body
  \item a declaration of several names
  \item a header with a macro
  \end{itemize}
\item[Again] hdl-vet skips the operation when the module has no such parameter.
\item[Sign-off] The sign-off requires \code{param-changed} at every path.
\item[Example] \file{retire.py}.
\end{entry}

\lstinputlisting[style=script]{examples/scripts/retire.py}

\lstinputlisting[style=script]{examples/scripts/remove_timer.py}

\lstinputlisting[style=console]{examples/console/remove.txt}

\file{retire.py} removes the top-level output \code{irq_o} from the design.
It removes the assignment and the port in one group. The signal \code{irq} stays.
Nothing reads it any more, so hdl-vet adds a note. The same script removes the parameter
\code{ID}, which nobody uses. The empty parameter list goes too.

hdl-vet refuses \file{remove_timer.py}. If the script removes the second timer, bit 1 of
\code{irq_o} of \code{periph_wrap} has no driver, and something reads it.

\section{Expectations by hand, and raw edits}
\label{sec:byhand}

Sometimes no operation does what you need. Then you edit the text yourself, and you
tell the sign-off what to expect. No operation checks raw edits. The
sign-off checks them, but it checks only what it can see. It can see the structure of
the design and the diagnostics (chapter~\ref{ch:signoff}).

\begin{entry}{\sig{Vet.splice}}
\item[Writes] The operation replaces the bytes from \code{start} up to \code{end} of the
  file \code{path} with \code{text} (a \code{bytes} value). Here \code{path} is a
  \emph{file}. The offsets are offsets into \code{vet.text(path)}. The file is a file of
  the design, for example the attribute \code{path} of a module handle.
\item[Refused] hdl-vet refuses the operation in these cases:
  \begin{itemize}
  \item a file that is not a file of the design, or is not editable (an added file is
    not editable)
  \item a range that overlaps an earlier splice
  \item a range outside the file
  \item an offset that is not an \code{int}, a negative offset, or a \code{start} after
    \code{end}
  \item a \code{text} or an \code{old} that is not \code{bytes}
  \item a range whose bytes differ from \code{old}, if you gave it
  \end{itemize}
  Always give \code{old}. Then the splice cannot go wrong when the text changes. Every
  refusal is a \code{VetError}. Its message starts with \code{splice} and the name of the
  file (chapter~\ref{ch:scripts}). The operation accepts the same range with the same text a
  second time, without a message.
\item[Again] A splice has no skip rule. With \code{old}, a second splice on text that
  is already edited is refused, because the bytes differ.
\item[Example] \file{by_hand.py}. After a splice, the next operation loads the design
  again, because hdl-vet does not know the effect of a splice.
\end{entry}

\begin{entry}{\sig{Vet.text}}
\item[Writes] The operation writes nothing. It returns the text of a file of the design,
  as \code{bytes}, as of the last load. The text does not have edits that are still in the
  queue, such as a raw splice, until the next load. A file that is not a file of the design
  is a \code{VetError}. So is a \code{path} that is neither a \code{str} nor a \code{Path}.
  Its message starts with \code{text} and the name of the file.
\end{entry}

\begin{entry}{\sig{Vet.expect}}
\item[Writes] The operation writes nothing in the text. It declares a change that the
  sign-off must see. \code{kind}, \code{path} and \code{name} are patterns (\code{*},
  \code{?} and \code{[...]}). Capital and small letters differ. \code{kind} is one of
  these:
  \begin{itemize}
  \item \code{instance-added}
  \item \code{instance-removed}
  \item \code{module-changed}
  \item \code{param-changed}
  \item \code{port-added}
  \item \code{port-removed}
  \item \code{port-changed}
  \item \code{port-type-changed}
  \item \code{connection-changed}
  \end{itemize}
  \code{path} is the hierarchical path of the instance. \code{name} is the port or
  parameter. It is empty for the kinds that name none. With \code{exact=True},
  \code{path} is a plain path and not a pattern.
\item[Rules] A change that no expectation or permission covers makes the sign-off fail
  (\emph{unexpected change}). An expectation that matches no change also makes the
  sign-off fail (\emph{expected change not seen}).
\item[Example] \file{by_hand.py}, \file{unexpected.py}, \file{unseen.py}.
\end{entry}

\begin{entry}{\sig{Vet.permit}}
\item[Writes] The operation writes nothing in the text. It works as
  \apiname{Vet.expect}, but the change \emph{can} happen. If no change happens, this is
  not an error. Use it for changes that follow from yours but that you do not have to see.
\item[Example] \file{by_hand.py}.
\end{entry}

\begin{entry}{\sig{Instance.allow_changes_below}}
\item[Writes] The operation writes nothing in the text. It permits every structural
  change strictly below each path of the instance.
\item[Example] \file{below.py}.
\end{entry}

\lstinputlisting[style=script]{examples/scripts/by_hand.py}

\lstinputlisting[style=script]{examples/scripts/below.py}

\lstinputlisting[style=script]{examples/scripts/unexpected.py}

\lstinputlisting[style=script]{examples/scripts/unseen.py}

\lstinputlisting[style=console]{examples/console/by_hand.txt}

\file{by_hand.py} makes two raw edits. It changes the tie-off of \code{test_i}, and
declares this with \code{expect}. It changes the parameter \code{STAGES}, and says
with \code{permit} that this can happen. \file{below.py} edits the enable of the
first timer. The change is below \code{u_periph}, and \code{allow_changes_below}
covers it.

The last two scripts show the two failures of the sign-off:

\begin{itemize}
\item a change that nobody declared
\item a declaration without a change
\end{itemize}
